\documentclass[11pt,final,hidelinks]{article}
\usepackage{lmodern}
\usepackage[T1]{fontenc}
\usepackage[english]{babel}
\usepackage[utf8]{inputenc}
\usepackage[activate={true,nocompatibility},final,tracking=true,kerning=true,
    spacing=true,factor=1100,stretch=10,shrink=10]{microtype}
\microtypecontext{spacing=nonfrench}
\usepackage[margin=1in]{geometry}
\usepackage{graphicx}
\graphicspath{{img/}}
\usepackage{caption}
\usepackage{subcaption}
\usepackage{booktabs}
\usepackage{array}
\usepackage{mathtools}
\usepackage{enumitem}
\usepackage{commath}
\usepackage{appendix}
\renewcommand\appendixtocname{Appendices}
\usepackage{amsmath}
\usepackage{amsthm}
\usepackage{amssymb}
\usepackage{pgfplots}
\usepackage{tikz}
\usepackage{tikzscale}
\usepackage[square,numbers]{natbib}
\bibliographystyle{plain}
\usepackage{siunitx}
\numberwithin{equation}{section}
\usepackage[hypertexnames=false]{hyperref}
\usepackage{cleveref}

\pgfplotsset{compat=newest}

\usepackage[fancy,section]{alex}

\newcommand{\Fone}{\RV{F_1}}
\newcommand{\fone}{\mathrm{f}_1}

\hypersetup{
    pdftitle={Binary classification measures for random approximate sets},
    pdfauthor={Alexander Towell},
    pdfsubject={computer science},
    pdfkeywords={
        probabilistic data structure,
        approximate set,
        binary classification,
        precision,
        recall,
        positive predictive value,
        F1-score,
        accuracy},
    colorlinks=true,
    linkcolor=magenta,
    citecolor=green,
    filecolor=blue,
    urlcolor=green
}

\title
{
    Binary classification measures\\
    \large for random approximate sets
}
\author
{
    Alexander Towell\\
    \texttt{atowell@siue.edu}
}
\date{2026}

\begin{document}
\maketitle
\begin{abstract}
We derive the probability distributions, expected values, and variances of standard binary classification measures (positive predictive value or precision, negative predictive value, accuracy, the $\fone$-score, and Youden's $J$ statistic) for the Bernoulli set model of random approximate sets with independent, element-wise errors.
The expected PPV and $\fone$-score are obtained via standard second-order delta method expansions whose accuracy we validate by Monte Carlo simulation.
All measures depend critically on the \emph{prevalence} (base rate): even a small false positive rate can drive precision to zero when the base rate is low, the classical predictive-value relation of diagnostic testing.
These are known statistical results restated for the model; no novelty is claimed.
The results apply exactly to any data structure satisfying the Bernoulli axioms, and approximately to Bloom filters and perfect hash filters, whose errors share a random build.
\end{abstract}

% Status and errata box (October 2026 freeze). \input right after \end{abstract}.
\begingroup
\setlength{\fboxsep}{7pt}%
\begin{center}
\fbox{\begin{minipage}{0.93\linewidth}\small
\setlength{\parskip}{4pt}%
\setlength{\parindent}{0pt}%
\textbf{Status and errata (October 2026).}
This is the final version of this preprint; no further revisions are planned.
It is superseded by Chapter~7 (\emph{Classification Measures}) of the monograph
\emph{Bernoulli Types: Two-Sided Error in Approximate Membership, Retrieval,
and Computation} (A.~Towell, 2026),
\url{https://metafunctor.com/latex/bernoulli-monograph/book.html}
(project hub: \url{https://metafunctor.com/research/bernoulli/}).
Please cite the monograph in preference to this preprint.

\textbf{Corrections in this version.}
\begin{itemize}\setlength{\itemsep}{1pt}\setlength{\parskip}{0pt}
\item Abstract, \cref{sec:intro}: the model is no longer called ``compositionally closed''; a two-rate pair does not propagate through arbitrary set expressions (see the companion composition paper).
\item Abstract, \cref{sec:conclusion_cm}: Bloom filters and perfect hash filters satisfy the axioms only approximately (errors share one random build); ``including Bloom filters'' is qualified.
\item \Cref{sec:prelim}: $\Expect{\beta}$ corrected from $\tprate$ to $\fnrate$ ($\beta$ is the realized false negative rate, consistent with $\mathrm{TPR}_p = 1 - \beta_p$).
\item \Cref{sec:interval_measures}: the interval PPV bounds hold for the ratio of expected counts, not for $\Expect{\PPV}$ itself.
\item \Cref{sec:intro}: novelty claims removed; the prevalence effect is the classical predictive-value relation.
\end{itemize}

\textbf{Predecessors credited.} The per-element channel is randomized
response~\cite{warner1965randomized}; misclassification in $2 \times 2$ tables is
treated by Bross~\cite{bross1954misclassification}; the dependence of predictive
values on prevalence is classical in diagnostic testing~\cite{altman1994predictive};
probabilistic precision, recall, and F-score are due to Goutte and
Gaussier~\cite{goutte2005probabilistic}; the expansions are the standard delta
method~\cite{casellaBerger}. No novelty is claimed.
A dated list of changes is in \texttt{ERRATA.md} in the source repository.
\end{minipage}}
\end{center}
\endgroup


\microtypesetup{protrusion=false}
\tableofcontents
\microtypesetup{protrusion=true}

%%%%%%%%%%%%%%%%%%%%%%%%%%%%%%%%%
% Introduction
%%%%%%%%%%%%%%%%%%%%%%%%%%%%%%%%%

\section{Introduction}
\label{sec:intro}

Binary classification measures---precision, recall, accuracy, F-score---are fundamental to evaluating information retrieval systems and probabilistic data structures.
When a data structure approximates a set, returning membership queries that may err, the standard measures become random variables whose distributions depend on the error model.
This paper derives those distributions within the Bernoulli set model~\cite{bernoulliSets}, a probabilistic framework for random approximate sets with independent, element-wise errors parameterized by false positive and false negative rates.

\paragraph{Motivation.}
In many applications the question is not ``what is the false positive rate?'' but rather ``given a positive test result, what is the probability it is correct?''
This is the positive predictive value (precision), and its answer depends critically on the prevalence (base rate) of positives.
A Bloom filter~\cite{bf} with a $1\%$ false positive rate achieves near-perfect precision when queried against a large positive set, but its precision collapses when positives are rare. This is the classical dependence of predictive values on prevalence in diagnostic testing~\cite{altman1994predictive}; we restate it for approximate data structures.

\paragraph{Scope.}
We collect exact distributions, closed-form approximations, and variances for PPV, NPV, accuracy, the $\fone$-score, and Youden's $J$ within the Bernoulli set model~\cite{bernoulliSets}.
The approximations are standard delta-method expansions~\cite{casellaBerger} applied to independent binomial counts, and no novelty is claimed for them.
The expected PPV and $\fone$-score are obtained via second-order delta method expansions (\cref{thm:approx_expected_precision,thm:f1_approx}), and their accuracy is validated by Monte Carlo simulation (\cref{fig:ppv_theory_vs_sim,fig:f1_theory_vs_sim}).
We show that all measures are functions of four parameters: the false positive rate $\fprate$, the true positive rate $\tprate$, the number of positives $\p$, and the number of negatives $\n$.

\paragraph{Related work.}
Fawcett~\cite{fawcettROC} provides an introduction to ROC analysis; Powers~\cite{powersEval} gives a systematic comparison of evaluation measures.
Goutte and Gaussier~\cite{goutte2005probabilistic} give a probabilistic treatment of precision, recall, and the F-score, with distributions inferred from observed counts; here the rates and the numbers of positives and negatives are known and the measures are sampling distributions.
The per-element error channel is that of randomized response~\cite{warner1965randomized}, and the effect of a known misclassification channel on $2 \times 2$ tables goes back to Bross~\cite{bross1954misclassification}.
The underlying Bernoulli set model and its compositional algebra are developed in a companion paper~\cite{bernoulliSets}; here we focus exclusively on the classification measures that the model induces.

\paragraph{Organization.}
\Cref{sec:prelim} restates the Bernoulli axioms and the binomial distributions of error counts.
\Cref{sec:measures} derives the distributions of PPV, NPV, accuracy, Youden's $J$, and the $\fone$-score.
\Cref{sec:conclusion_cm} concludes.
Appendix~\ref{sec:proof_approx_expected_precision} gives the full delta-method proof for expected PPV, Appendix~\ref{sec:proof_f1} gives the proof for the $\fone$-score, and Appendix~\ref{app:samp} develops the sampling distribution framework.

%%%%%%%%%%%%%%%%%%%%%%%%%%%%%%%%%
% Preliminaries
%%%%%%%%%%%%%%%%%%%%%%%%%%%%%%%%%

\section{Preliminaries}
\label{sec:prelim}

We briefly restate the elements of the Bernoulli set model~\cite{bernoulliSets} needed in this paper.

\begin{axiom}[Element-wise independence]
\label{asm:element_indep}
For all distinct $x, y \in U$, the error events at $x$ and $y$ are mutually independent:
\begin{equation}
P(\mathbf{1}_{\tilde{A}}(x) \neq \mathbf{1}_{A}(x) \mid \mathbf{1}_{\tilde{A}}(y) \neq \mathbf{1}_{A}(y)) =
	P(\mathbf{1}_{\tilde{A}}(x) \neq \mathbf{1}_{A}(x)).
\end{equation}
More generally, the family $\{\mathbf{1}_{\tilde{A}}(x) : x \in U\}$ is mutually independent given the objective set $A$ and the error rates.
Furthermore, within each partition block $B_i$, the error indicators $\{\mathbf{1}_{\tilde{A}}(x) : x \in B_i\}$ are identically distributed.
\end{axiom}

\begin{proposition}
\label{prop:element_rates}
By \cref{asm:element_indep}, for any single element $x \notin A$,
$\Prob{\mathbf{1}_{\tilde{A}}(x) = 1} = \fprate$, and for $x \in A$,
$\Prob{\mathbf{1}_{\tilde{A}}(x) = 1} = \tprate$.
The expected sample rates satisfy $\Expect{\alpha} = \fprate$ and $\Expect{\beta} = \fnrate$, where $\alpha$ and $\beta$ are the realized false positive and false negative rates.
\end{proposition}
\begin{proof}
See~\cite{bernoulliSets}.
\end{proof}

\begin{theorem}
\label{thm:fpbinom}
The random number of false positives $\FP$ given $\RV{N} = n$ in the second-order random approximate set $\tilde{\RV{R}}[\fprate]$ is given by
\begin{equation}
	\FP_n = n \alpha
\end{equation}
with a distribution given by
\begin{equation}
    \FP_n \sim \bindist(n, \fprate).
\end{equation}
\end{theorem}
\begin{proof}
See~\cite{bernoulliSets}.
\end{proof}

\begin{theorem}
\label{thm:fpr}
The random false positive rate $\alpha$ conditioned on $\RV{N} = n$ is denoted by $\alpha_n$ and has a distribution given by
\begin{equation}
    \alpha_n = \frac{\FP_n}{n},
\end{equation}
with an expectation $\fprate$, variance $\fprate(1-\fprate) / n$, and probability mass function
\begin{equation}
	p_{\alpha_n}(\fprateob | \fprate) = p_{\FP_n}(\fprateob n | \fprate)
\end{equation}
over the support $\SetBuilder{ \frac{j}{n} \in \RatSet}{j \in \{0,\ldots,n\}}$.
\end{theorem}
\begin{proof}
See~\cite{bernoulliSets}.
\end{proof}

\begin{corollary}
\label{cor:tnbinom}
Given $n$ negatives, the number of \emph{true negatives} in a random approximate set with a false positive rate $\fprate$ is a random variable denoted by $\TN_n$ with a distribution given by
\begin{equation}
	\TN_n = n - \FP_n \sim \bindist\!\left(n, 1-\fprate\right).
\end{equation}
By definition, the \emph{true negative rate} $\mathrm{TNR}_n = \TN_n / n = 1 - \alpha_n$.
\end{corollary}
\begin{proof}
See~\cite{bernoulliSets}.
\end{proof}

\begin{theorem}
\label{thm:fnbinom}
Given $p$ positives, the uncertain number of \emph{false negatives} in random approximate sets with a false negative rate $\fnrate$ is modeled as a binomial distributed random variable denoted by $\FN_p$,
\begin{equation}
	\FN_p \sim \bindist(p, \fnrate).
\end{equation}
\end{theorem}
\begin{proof}
See~\cite{bernoulliSets}.
\end{proof}

\begin{corollary}
\label{cor:tpbinom}
Given $p$ positives, the number of \emph{true positives} in an approximate set
with a false negative rate $\fnrate$ is a random variable denoted by $\TP_p$
with a distribution given by
\begin{equation}
    \TP_p \sim \bindist(p, \tprate).
\end{equation}
By definition, the \emph{true positive rate} is given by $\TPR_p = 1 -
\beta_p$.
\end{corollary}
\begin{proof}
See~\cite{bernoulliSets}.
\end{proof}

%%%%%%%%%%%%%%%%%%%%%%%%%%%%%%%%%
% Binary classification measures
%%%%%%%%%%%%%%%%%%%%%%%%%%%%%%%%%

\section{Binary classification measures}
\label{sec:measures}

When a random approximate set $\ASet{A}$ replaces an objective set $\Set{A}$, any function of the membership query results becomes a random variable whose distribution is inherited from the binomial error counts of \cref{sec:prelim}.
Binary classification measures---positive predictive value, negative predictive value, accuracy, the $\fone$-score, and Youden's $J$---are such functions.
These measures are standard in binary classification and information retrieval; see Fawcett~\cite{fawcettROC} for an introduction to ROC analysis and Powers~\cite{powersEval} for a systematic comparison of evaluation measures.

We derive expected values and variances for each measure.
Since PPV, NPV, and the $\fone$-score are nonlinear ratios of binomial random variables, their expectations require the delta method (a Taylor expansion about the means); accuracy and Youden's $J$, being linear combinations, admit exact closed-form moments.
The accuracy of \emph{predictions} about objective sets given a corresponding approximate set is usually the more relevant performance measure.
The \emph{positive predictive value} is given by the following definition.
\begin{definition}
	The \emph{positive predictive value} is a performance measure defined as
	\begin{equation}
	\ppv = \frac{t_p}{t_p + f_p}
	\end{equation}
	where $t_p$ is the number of \emph{true positives} and $f_p$ is the number of \emph{false positives}.
\end{definition}

The positive predictive value of random approximate sets is a random variable given by the following theorem.
\begin{theorem}
	\label{thm:approx_expected_precision}
	Given $\n$ negatives, $\p$ positives, and a random approximate set with false positive and true positive rates $\fprate$ and $\tprate$ respectively, the \emph{positive predictive value} is a random variable
	\begin{equation}
	\PPV = \frac{\TP_{\p}}{\TP_{\p} + \FP_{\n}}
	\end{equation}
	with an \emph{expectation} given \emph{approximately} by the second-order delta method~\cite{casellaBerger} (a Taylor expansion of $\Expect{\TP_\p / (\TP_\p + \FP_\n)}$ about the means, valid when $\p\tprate$ and $\n\fprate$ are both large; the error is $\mathcal{O}(1/\min(\p,\n)^2)$)
	\begin{equation}
	\label{eq:approx_expected_precision}
	\ppv(\tprate, \fprate, \p, \n) \approx
	\frac{\overline{t}_\p}{\overline{t}_\p + \overline{f}_\p} +
	\frac{\overline{t}_\p \sigma_{\!f_\p}^2 - \overline{f}_\p
		\sigma_{\!t_\p}^2}{\left(\overline{t}_\p + \overline{f}_\p\right)^3},
	\end{equation}
	where $\overline{t}_\p = \p \tprate$ is the expected \emph{true positive} frequency, $\overline{f}_\p =  \n \fprate$ is the expected \emph{false positive} frequency, $\sigma_{\!t_\p}^2 = (1-\tprate) \overline{t}_\p$ is the variance of the \emph{true positive} frequency, and $\sigma_{\!f_\p}^2 = (1-\fprate) \overline{f}_\p$ is the variance of \emph{false positive} frequency.
	The first-order delta method variance is
	\begin{equation}
	\label{eq:var_ppv}
	\Var{\PPV} \approx
	\frac{\overline{f}_\p^{\,2}\,\sigma_{\!t_\p}^2 +
		\overline{t}_\p^{\,2}\,\sigma_{\!f_\p}^2}
	{\left(\overline{t}_\p + \overline{f}_\p\right)^4}.
	\end{equation}
\end{theorem}
See \cref{sec:proof_approx_expected_precision} for a proof of
\cref{thm:approx_expected_precision}.

\begin{figure}
	\def\svgwidth{\columnwidth/2}
	%\def\svgscale{0.75}
	\centering
	\captionsetup{justification=centering}
	\input{img/out.pdf_tex}
	\caption{Relative frequency of positive predictive values for several different parameterizations of the false positive and true positive rates given $\n = 900$ negatives and $\p=100$ positives.}
	\label{fig:prec_vs_fprate_and_fnrate}
\end{figure}

\begin{figure}[ht]
	\centering
	\begin{tikzpicture}
	\begin{axis}[
		width=0.7\textwidth,
		height=0.45\textwidth,
		xlabel={False positive rate $\fprate$},
		ylabel={Expected precision (PPV)},
		xmin=0, xmax=0.05,
		ymin=0, ymax=1,
		legend pos=north east,
		ymajorgrids, xmajorgrids,
		major grid style={gray!30},
		thick,
	]
	\addplot[blue, solid] table[col sep=comma, x index=0, y index=1]
		{data/prec_vs_fprate3.csv};
	\addlegendentry{$\lambda = 10\%$}
	\addplot[red, dashed] table[col sep=comma, x index=0, y index=1]
		{data/prec_vs_fprate2.csv};
	\addlegendentry{$\lambda = 1\%$}
	\addplot[black!60!green, dashdotted] table[col sep=comma, x index=0, y index=1]
		{data/prec_vs_fprate1.csv};
	\addlegendentry{$\lambda = 0.1\%$}
	\end{axis}
	\end{tikzpicture}
	\caption{Expected precision (positive predictive value) as a function
	of the false positive rate $\fprate$ for a positive Bernoulli set
	($\fnrate = 0$) at three prevalences
	$\lambda = \p/(\p+\n)$.
	Even a small false positive rate dramatically reduces precision
	when the base rate is low.}
	\label{fig:ppv_vs_fprate}
\end{figure}

\begin{figure}[ht]
	\centering
	\begin{tikzpicture}
	\begin{axis}[
		width=0.7\textwidth,
		height=0.45\textwidth,
		xlabel={False positive rate $\fprate$},
		ylabel={Expected precision (PPV)},
		xmin=0, xmax=0.05,
		ymin=0, ymax=1,
		legend pos=north east,
		ymajorgrids, xmajorgrids,
		major grid style={gray!30},
		thick,
	]
	\addplot[red, dashed] table[col sep=comma, x index=0, y index=1]
		{data/prec_vs_fprate2.csv};
	\addlegendentry{Theory ($\lambda = 1\%$)}
	\addplot[only marks, mark=o, mark size=1.2pt, blue]
		table[col sep=comma, x index=0, y index=1]
		{data/ppv_sim_vs_theory.csv};
	\addlegendentry{Monte Carlo ($N = 50\,000$)}
	\end{axis}
	\end{tikzpicture}
	\caption{Comparison of theoretical expected PPV
	(\cref{eq:approx_expected_precision}, dashed) with Monte Carlo simulation
	(circles, $N = 50\,000$ trials per point) for a positive Bernoulli set
	($\fnrate = 0$) with prevalence $\lambda = 1\%$ ($\p = 100$, $\n = 9900$).
	The theoretical curve and simulation points are indistinguishable,
	validating the second-order Taylor approximation.}
	\label{fig:ppv_theory_vs_sim}
\end{figure}

We make the following observations about \cref{eq:approx_expected_precision}:
\begin{enumerate}
	\item For sufficiently large approximate sets, $\ppv \approx \overline{t}_\p / (\overline{t}_\p + \overline{f}_\p)$.
	\item If $\fprate \neq 0$, as $\n \to \infty$, $\ppv \to 0$.
	\item As $\fprate \to 0, \ppv \to 1$.
\end{enumerate}

\emph{Accuracy} is given by the following definition.
\begin{definition}
	The \emph{accuracy} is the proportion of true results (both \emph{true
		positives} and \emph{true negatives}) in the universe of positives and
	negatives, $(t_p + t_n)/(\p + \n)$, where $t_p$, $t_n$, $\p$, and $\n$ are
	respectively the number of \emph{true positives}, \emph{true negatives},
	\emph{positives}, and \emph{negatives}.
\end{definition}

The \emph{expected} accuracy is given by the following theorem.
Let $\lambda = \p / (\p + \n)$ denote the \emph{base rate} (proportion of positives).
\begin{theorem}
	\label{thm:approx_expected_accuracy}
	Given $\p$ positives and $\n$ negatives, the accuracy of a random approximate set with an
	\emph{expected} false positive rate $\fprate$ and an \emph{expected} true
	positive rate $\tprate$ is a random variable given by
	\begin{equation}
	\ACC_{\p + \n} = \lambda \mathrm{TPR}_\p + (1 - \lambda) \mathrm{TNR}_n
	\end{equation}
	with an \emph{expected} accuracy
	\begin{equation}
	\label{eq:approx_expected_accuracy}
	\acc(\tprate,\fprate,\n,\p) = \lambda \tprate + (1-\lambda) \tnrate
	\end{equation}
	and a variance
	\begin{equation}
	\frac{\lambda \fnrate \tprate + (1 - \lambda)\fprate
		\tnrate}{\p+\n}.
	\end{equation}
\end{theorem}
\begin{proof}
	Suppose the $u$ elements in the universe can be partitioned into $\p$
	positives and $\n$ negatives. An approximate set $\tilde{S}$ with a false
	positive rate $\fprate$ and false negative rate $\fnrate$ has an uncertain
	accuracy
	\begin{equation}
	\ACC_{\p + \n} = \frac{\TP_\p + \TN_\n}{\p + \n}.
	\end{equation}
	The expected accuracy is given by the expectation
	\begin{align}
	\Expect{\ACC_{\p + \n}}
	&= \Expect{\frac{\TP_\p + \TN_\n}{\p + \n}}\\
	&= \frac{\p (1 - \fnrate) + \n (1 - \fprate)}{\p + \n}.
	\end{align}
	Noting that $\n/(\p+\n) = 1 - \p/(\p+\n)$ and letting $\lambda = \p/(\p +
	\n)$,
	\begin{equation}
	\Expect{\ACC_{\p + \n}} = \lambda (1 - \fnrate) + (1 - \lambda) (1 -
	\fprate).
	\end{equation}
	The variance
	\begin{align}
	\Var{\ACC_{\p + \n}}
	&= \Var{\frac{\TP_\p}{\p + \n}} + \Var{\frac{\TN_\n}{\p + \n}}\\
	&= \frac{1}{(\p + \n)^2}\Var{\TP_\p} + \frac{1}{(\p + \n)^2}\Var{\TN_\n}\\
	&= \frac{\p \fnrate (1 - \fnrate)}{(\p + \n)^2} + \frac{\n \fprate (1 -
		\fprate)}{(\p + \n)^2}\\
	&= \frac{\lambda \fnrate \tprate + (1 - \lambda)\fprate
		\tnrate}{\p+\n}.
	\end{align}
\end{proof}

\begin{definition}
	The \emph{negative predictive value} is a performance measure defined as
	\begin{equation}
	\npv = \frac{t_n}{t_n + f_n}
	\end{equation}
	where $t_n$ and $f_n$ are respectively the number of \emph{true negatives}
	and \emph{false negatives}.
\end{definition}
The expected negative predictive value is given by the following theorem.
\begin{theorem}
	\label{thm:npv_approx}
	Given $\p$ positives, $\n$ negatives, and a random approximate set with false positive and true positive rates $\fprate$ and $\tprate$ respectively, the negative predictive value is a \emph{random variable}
	\begin{equation}
	\NPV = \frac{\TN_{\n}}{\TN_{\n} + \FN_{\p}}
	\end{equation}
	with an \emph{expectation} given approximately by
	\begin{equation}
	\label{eq:npv_approx}
	\npv(\tprate, \fprate , \p, \n) \approx
	\frac{\overline{t}_n}{\overline{t}_n + \overline{f}_{\!n}} +
	\frac{\overline{t}_n \sigma_{\!f_n}^2 - \overline{f}_{\!n}
		\sigma_{\!t_n}^2}{\left(\overline{t}_n + \overline{f}_{\!n}\right)^3},
	\end{equation}
	where $\overline{t}_\n = \n(1-\fprate)$ is the expected \emph{true negative} frequency, $\overline{f}_{\!\n} =  \p(1-\tprate)$ is the expected \emph{false negative} frequency, $\sigma_{\!t_\n}^2 = \fprate \overline{t}_\n$ is the variance of the \emph{true negative} frequency, and $\sigma_{\!f_\n}^2 = \tprate\overline{f}_{\!\n}$ is the variance of the \emph{false negative} frequency.
	The first-order delta method variance is
	\begin{equation}
	\label{eq:var_npv}
	\Var{\NPV} \approx
	\frac{\overline{f}_{\!\n}^{\,2}\,\sigma_{\!t_\n}^2 +
		\overline{t}_\n^{\,2}\,\sigma_{\!f_\n}^2}
	{\left(\overline{t}_\n + \overline{f}_{\!\n}\right)^4}.
	\end{equation}
\end{theorem}
The proof for \cref{thm:npv_approx} follows the same pattern as the proof for \cref{thm:approx_expected_precision}.

Youden's $J$ statistic is a measure of the performance of a binary test, defined as
\begin{equation}
J = \frac{t_p}{t_p + f_n} + \frac{t_n}{t_n + f_p} - 1,
\end{equation}
with a range $[-1,1]$.
In the case of the random approximate set model, $J$ is a random variable
\begin{equation}
\RV{J} = \mathrm{TPR}_\p - \alpha_n,
\end{equation}
which has an \emph{expectation}
\begin{equation}
\Expect{\RV{J}} = \tprate - \fprate
\end{equation}
and, by the independence of $\mathrm{TPR}_\p$ and $\alpha_n$, a \emph{variance}
\begin{equation}
\label{eq:var_j}
\Var{\RV{J}} = \frac{\tprate\,\fnrate}{\p} + \frac{\fprate\,\tnrate}{\n}.
\end{equation}

\subsection{The \texorpdfstring{$\fone$}{F1}-score}
\label{sec:f1}
\begin{definition}
	The \emph{$\fone$-score} is the harmonic mean of precision and recall:
	\begin{equation}
	\fone = \frac{2\,t_p}{2\,t_p + f_p + f_n}.
	\end{equation}
\end{definition}
Since $f_n = \p - t_p$, the $\fone$-score may be written as a function of $t_p$ and $f_p$ alone:
\begin{equation}
\label{eq:f1_reduced}
\fone = \frac{2\,t_p}{t_p + f_p + \p}.
\end{equation}
\begin{theorem}
	\label{thm:f1_approx}
	Given $\p$ positives, $\n$ negatives, and a random approximate set with false positive and true positive rates $\fprate$ and $\tprate$ respectively, the $\fone$-score is a random variable
	\begin{equation}
	\Fone = \frac{2\,\TP_{\p}}{\TP_{\p} + \FP_{\n} + \p}
	\end{equation}
	with an \emph{expectation} given approximately by the second-order delta method
	\begin{equation}
	\label{eq:f1_approx}
	\Expect{\Fone} \approx
	\frac{2\,\overline{t}_\p}{\overline{t}_\p + \overline{f}_\p + \p}
	+ \frac{2\,\overline{t}_\p\,\sigma_{\!f_\p}^2 -
		2\left(\overline{f}_\p + \p\right)\sigma_{\!t_\p}^2}
	{\left(\overline{t}_\p + \overline{f}_\p + \p\right)^3}
	\end{equation}
	and a first-order delta method \emph{variance}
	\begin{equation}
	\label{eq:var_f1}
	\Var{\Fone} \approx
	\frac{4\left(\overline{f}_\p + \p\right)^2 \sigma_{\!t_\p}^2 +
		4\,\overline{t}_\p^{\,2}\,\sigma_{\!f_\p}^2}
	{\left(\overline{t}_\p + \overline{f}_\p + \p\right)^4},
	\end{equation}
	where $\overline{t}_\p = \p\tprate$, $\overline{f}_\p = \n\fprate$,
	$\sigma_{\!t_\p}^2 = \p\tprate(1-\tprate)$, and
	$\sigma_{\!f_\p}^2 = \n\fprate(1-\fprate)$.
\end{theorem}
See \cref{sec:proof_f1} for a proof of \cref{thm:f1_approx}.

\begin{figure}[ht]
	\centering
	\begin{tikzpicture}
	\begin{axis}[
		width=0.7\textwidth,
		height=0.45\textwidth,
		xlabel={False positive rate $\fprate$},
		ylabel={Expected $\fone$-score},
		xmin=0, xmax=0.05,
		ymin=0, ymax=1,
		legend pos=north east,
		ymajorgrids, xmajorgrids,
		major grid style={gray!30},
		thick,
	]
	\addplot[red, dashed] table[col sep=comma, x index=0, y index=1]
		{data/f1_theory.csv};
	\addlegendentry{Theory ($\lambda = 1\%$)}
	\addplot[only marks, mark=o, mark size=1.2pt, blue]
		table[col sep=comma, x index=0, y index=1]
		{data/f1_sim_vs_theory.csv};
	\addlegendentry{Monte Carlo ($N = 50\,000$)}
	\end{axis}
	\end{tikzpicture}
	\caption{Comparison of theoretical expected $\fone$-score
	(\cref{eq:f1_approx}, dashed) with Monte Carlo simulation
	(circles, $N = 50\,000$ trials per point) for a positive Bernoulli set
	($\fnrate = 0$) with prevalence $\lambda = 1\%$ ($\p = 100$, $\n = 9900$).
	The theoretical curve and simulation points are indistinguishable,
	validating the second-order delta method approximation.}
	\label{fig:f1_theory_vs_sim}
\end{figure}

\begin{table}
	\centering
	\caption{Various \emph{expected} performance measures.}
	\label{tbl:perf_sum}
	\begin{tabular}{@{} l l l @{}}
		\toprule
		\textbf{measure} & \textbf{parameter} & \textbf{expected value}\\
		\midrule
		true positive rate & $\tpr(\tprate)$ & $\tprate$\\
		false positive rate & $\fpr(\fprate)$ & $\fprate$\\
		false negative rate & $\fnr(\tprate)$ & $1-\tprate$\\
		true negative rate & $\tnr(\fprate)$ & $1-\fprate$\\
		accuracy & $\acc$ & \cref{eq:approx_expected_accuracy}\\
		positive predictive value & $\ppv$ & \cref{eq:approx_expected_precision}\\
		negative predictive value & $\npv$ & \cref{eq:npv_approx}\\
		$\fone$-score & $\fone$ & \cref{eq:f1_approx}\\
		false discovery rate & $\fdr$ & $1-\ppv$\\
		false omission rate & $\forFn$ & $1-\npv$\\
		\bottomrule
	\end{tabular}
\end{table}
\Cref{tbl:perf_sum} may be used to reparameterize an approximate set.
\begin{example}
	Suppose we seek a \emph{positive approximate set} with an expected accuracy
	$\gamma$. By \cref{tbl:perf_sum},
	\begin{equation}
	\gamma = \acc(\fprate,\fnrate=0,\lambda) = 1 - \fprate(1 - \lambda).
	\end{equation}
	% Y = 1 - e(1- lam)
	% (1-Y)/(1-lam) = e
	Solving for $\fprate$ in terms of $\gamma$ yields the result
	\begin{equation}
	\fprate(\gamma, \lambda) = \frac{1 - \gamma}{1 -
		\lambda}\
	\end{equation}
	subject to $0 \leq \lambda \leq \gamma \leq 1$ and $\lambda < 1$.
	Under this parameterization of the 	positive approximate set, $\lambda$
	must be known (or estimated).
	Note that if $\lambda = 1$ then $\fprate(\gamma, \lambda=1)$ is undefined as expected, but as $\lambda$ goes to $1$, $\fprate(\,\cdot\,; \lambda)$ goes to $1$ and $\gamma$ goes to $1$, which logically follows since if there are no negatives, there can be no false positives.
\end{example}

\subsection{Interval-valued classification measures}
\label{sec:interval_measures}
When the error rates or the proportion of positives are not known exactly, interval arithmetic~\cite{bernoulliComposition,basicinterval,mooreInterval} may be applied: each uncertain parameter is replaced by an interval guaranteed to contain its true value, and the classification-measure formulas are evaluated over these intervals to produce guaranteed bounds.

\begin{example}[Interval accuracy]
\label{ex:interval_accuracy}
	Suppose we wish to determine the expected accuracy given that the proportion of positives $\lambda$ is known to be some value in the interval $[\lambda]$, and that the false positive and false negative rates lie in $[\fprate]$ and $[\fnrate]$ respectively.
	The accuracy formula $\acc = \lambda(1-\fnrate) + (1-\lambda)(1-\fprate)$ may be rewritten as
	\begin{equation}
		\acc = (1-\fprate) + \lambda(\fprate - \fnrate).
	\end{equation}
	Since the accuracy is multilinear in $(\fprate, \fnrate, \lambda)$, its extremes over the box $[\fprate] \times [\fnrate] \times [\lambda]$ are attained at vertices.
	The formula is decreasing in $\fprate$ and $\fnrate$, but its monotonicity in $\lambda$ depends on the sign of $\fprate - \fnrate$:
	\begin{itemize}
		\item If $\fprate \geq \fnrate$ (the common case, e.g., Bloom filters), accuracy increases in $\lambda$, so:
		\begin{equation}
			\acc \in \bigl[
			(1-\overline{\fprate}) + \underline{\lambda}(\overline{\fprate} - \overline{\fnrate}),\;
			(1-\underline{\fprate}) + \overline{\lambda}(\underline{\fprate} - \underline{\fnrate})
			\bigr].
		\end{equation}
		\item If $\fprate < \fnrate$, accuracy decreases in $\lambda$, so the $\lambda$ endpoints are reversed:
		\begin{equation}
			\acc \in \bigl[
			(1-\overline{\fprate}) + \overline{\lambda}(\overline{\fprate} - \overline{\fnrate}),\;
			(1-\underline{\fprate}) + \underline{\lambda}(\underline{\fprate} - \underline{\fnrate})
			\bigr].
		\end{equation}
	\end{itemize}
	In practice, guaranteed bounds are obtained by evaluating the accuracy at all $2^3 = 8$ vertices and taking the minimum and maximum.
\end{example}

\begin{example}[Interval PPV]
\label{ex:interval_ppv}
Given interval-valued rates, the positive predictive value inherits interval bounds.
Since $\ppv \approx \overline{t}/(\overline{t} + \overline{f})$ for large sets (\cref{thm:approx_expected_precision}), substituting $\overline{t} = p\,\tprate$ and $\overline{f} = n\,\fprate$ and evaluating at the extremes of $[\fprate]$ and $[\tprate]$ yields
\begin{equation}
[\ppv] \approx \biggl[\frac{p\,\underline{\tprate}}{p\,\underline{\tprate} + n\,\overline{\fprate}},\; \frac{p\,\overline{\tprate}}{p\,\overline{\tprate} + n\,\underline{\fprate}}\biggr],
\end{equation}
These are guaranteed bounds on the ratio of expected counts $p\tprate/(p\tprate + n\fprate)$, the large-set approximation of the expected precision; they bound $\Expect{\PPV}$ itself only up to the delta-method correction of \cref{eq:approx_expected_precision}.
\end{example}

%%%%%%%%%%%%%%%%%%%%%%%%%%%%%%%%%
% Conclusion
%%%%%%%%%%%%%%%%%%%%%%%%%%%%%%%%%

\section{Conclusion}
\label{sec:conclusion_cm}

We have derived the probability distributions, closed-form expected values, and variances of the principal binary classification measures---positive predictive value, negative predictive value, accuracy, the $\fone$-score, and Youden's $J$ statistic---for the Bernoulli set model.
In each case the measure is a random variable whose distribution follows from the binomial error counts established by the element-wise independence axiom (\cref{asm:element_indep}).
The expected PPV and $\fone$-score were obtained via second-order delta method expansions (\cref{thm:approx_expected_precision,thm:f1_approx}), and their accuracy was validated by Monte Carlo simulation (\cref{fig:ppv_theory_vs_sim,fig:f1_theory_vs_sim}).

These results apply exactly to any data structure satisfying the Bernoulli axioms. Bloom filters~\cite{bf} and perfect hash filters~\cite{phf} satisfy them only approximately: errors at distinct elements share the random state of one build, so the realized rates vary from build to build and the classical Bloom filter rate is itself an approximation~\cite{bose2008fpr}.
In particular, the sensitivity of PPV to prevalence (\cref{fig:ppv_vs_fprate}) has practical implications for system design: a data structure with a nominally small false positive rate may still produce unreliable positive predictions when the base rate of positives is low.
Designers should therefore evaluate their systems not only by FPR and FNR but also by the predictive values that the intended query workload will induce.

Several directions remain open.
First, the distribution of the Matthews correlation coefficient (MCC) under the Bernoulli model has not been derived in closed form; it involves a more complex ratio of binomial random variables than PPV or $\fone$.
Second, extension to correlated error models---where the element-wise independence axiom is relaxed---would broaden the applicability of the framework to data structures with dependent hash functions or adversarial inputs.
Finally, the interplay between prevalence and the higher-order composition algebra of~\cite{bernoulliSets,bernoulliComposition} deserves further study: as approximate sets are composed through union and intersection, the effective prevalence at each stage changes, and so do the classification measures.


\appendices
%%%%%%%%%%%%%%%%%%%%%%%%%%%%%%%%%%
% Appendix
%%%%%%%%%%%%%%%%%%%%%%%%%%%%%%%%%%

\section{Proof of theorem~\ref{thm:approx_expected_precision}}
\label{sec:proof_approx_expected_precision}
Given $p$ positives and $n$ negatives, by \cref{thm:approx_expected_precision} an approximate set with a false positive rate $\fprate$ and a false negative rate $\fnrate$ has an \emph{expected} precision given \emph{approximately} by
\begin{equation*}
    \ppv(\fnrate, \fprate ; n, p) \approx \frac{\overline{t}}{\overline{t} + \overline{f}} +
    \frac{\overline{t} \sigma_{\!f}^2 - \overline{f} \sigma_{\!t}^2}{\left(\overline{t} + \overline{f}\right)^3},
    \tag{\ref{eq:approx_expected_precision} revisited}
\end{equation*}
where $\overline{t} = p \tprate$ is the \emph{expected} number of \emph{true positives}, $\overline{f} =  n \fprate$ is the \emph{expected} number of \emph{false positives}, $\sigma_{\!t}^2 = p \fnrate \tprate$ is the variance of the number of \emph{true positives}, and $\sigma_{\!f}^2 = n \fprate \tnrate$ is the variance of the number of \emph{false positives}.
\begin{proof}
The positive predictive value is a random variable given by
\begin{equation}
    \frac{\TP_p}{\TP_p + \FP_n}.
\end{equation}
We are interested in the \emph{expected} positive predictive value,
\begin{equation}
    \ppv(\fprate,\tprate) = \Expect{\frac{\TP_p}{\TP_p + \FP_n}}.
\end{equation}
This expectation is of a non-linear function of random variables, which is problematic so we choose to approximate the expectation.

Let the \emph{positive predictive value} function be denoted by
\begin{equation}
    f(t, f) = \frac{t}{t + f},
\end{equation}
where $t$ is the number of true positives and $f$ is the number of false positives.
We approximate this function with a second-order Taylor series.
The gradient of $f$ is given by
\begin{equation}
    \nabla f(t,f) =
    \frac{1}{(t + f)^2}
    \begin{bmatrix}
        f\\
        -t\\
    \end{bmatrix}
\end{equation}
and the Hessian of $f$ is given by
\begin{equation}
    \mathcal{H}(t,f) =
    \frac{1}{(t + f)^3}
    \begin{bmatrix}
        -2 f & t-f\\
        t-f & 2 t \\
    \end{bmatrix}
\end{equation}

A second-order Taylor approximation $\Fun{g}$ of $f$ about the expected value of $\TP_p$, denoted by $\overline{t}$, and the expected value of $\FP_n$, denoted by $\overline{f}$, is given by
\begin{equation}
    \Fun{g}(t,f) =
    f\left(\overline{t},\overline{f}\right) + \nabla f(\overline{t},\overline{f})^{\intercal}
    \begin{bmatrix}
        t - \overline{t}\\
        f - \overline{f}\\
    \end{bmatrix}
    + \frac{1}{2}
    \begin{bmatrix}
        t - \overline{t}\\
        f - \overline{f}\\
    \end{bmatrix}^{\intercal}
    \mathcal{H}(\overline{t},\overline{f})
    \begin{bmatrix}
        t - \overline{t}\\
        f - \overline{f}\\
    \end{bmatrix}
\end{equation}
As a function of random variables $\TP_p$ and $\FP_n$, $\Fun{g}\!\left(\TP_p,\FP_n\right)$ is a random variable.
Since $\Expect{\TP_p - \overline{t}} = 0$ and $\Expect{\FP_n - \overline{f}} = 0$, we immediately simplify the expectation of $\Fun{g}$ to
\begin{equation}
\label{eq:proof_hess1}
    \Expect{\Fun{g}(\TP_p,\FP_n)} = \frac{\overline{t}}{\overline{t}+\overline{f}} + \frac{\Expect{A}}{(\overline{t} + \overline{f})^3}
\end{equation}
where
\begin{equation}
    A = \frac{1}{2}
    \begin{bmatrix}
        \TP_p - \overline{t}\\
        \FP_n - \overline{f}
    \end{bmatrix}^{\intercal}
    \begin{bmatrix}
        -2 \overline{f} \left(\TP_p - \overline{t}\right) + \left(\overline{t}-\overline{f}\right)\left(\FP_n - \overline{f}\right)\\
        \left(\overline{t}-\overline{f}\right)\left(\TP_p - \overline{t}\right) + 2 \overline{t}\left(\FP_n - \overline{f}\right)
    \end{bmatrix}
\end{equation}
Expanding the quadratic form yields
\begin{equation}
    A = -\overline{f} \left(\TP_p - \overline{t}\right)^2 +
    \left(\overline{t}-\overline{f}\right)\left(\TP_p -
    \overline{t}\right)\left(\FP_n - \overline{f}\right) +
    \overline{t}\left(\FP_n - \overline{f}\right)^2
\end{equation}
As a linear operator, the expectation of $A$ is equivalent to
\begin{equation}
    \Expect{A} = -\overline{f} \Expect{(\TP_p - \overline{t})^2} + \left(\overline{t}-\overline{f}\right)\Expect{\left(\FP_n - \overline{f}\right)\left(\TP_p - \overline{t}\right)} + \overline{t}\Expect{(\FP_n - \overline{f})^2}
\end{equation}
By definition, $\Expect{(\TP_p - \overline{t})^2}$ is the variance of $\TP_p$, $\Expect{(\FP_n - \overline{f})^2}$ is the variance of $\FP_n$, and $\Expect{\left(\FP_n - \overline{f}\right)\left(\TP_p - \overline{t}\right)}$ is the covariance of $\TP_p$ and $\FP_n$, which is $0$ since they are independent.
Making these substitutions results in
\begin{equation}
    \Expect{A} = \overline{t} \Var{\FP_n} - \overline{f} \Var{\TP_p}
\end{equation}
Substituting this result into \cref{eq:proof_hess1} yields
\begin{equation}
    \Expect{\Fun{g}(\TP_p,\FP_n)} =
    \frac{\overline{t}}{\overline{t}+\overline{f}} +
    \frac{-\overline{f} \Var{\TP_p} +
    \overline{t}\Var{\FP_n}}{(\overline{t} + \overline{f})^3}
\end{equation}
By \cref{thm:fpbinom}, $\FP_n$ is binomially distributed with a mean $n \fprate$ and a variance $n \fprate \tnrate$ and by \cref{cor:tpbinom}, $\TP_p$ is binomially distributed with a mean $p \tprate$ and a variance $p \fnrate \tprate$.
Substituting $\Var{\FP_n} = n \fprate(1-\fprate) = \overline{f}(1-\fprate) = \sigma_f^2$ and $\Var{\TP_p} = p \tprate(1-\tprate) = \overline{t}(1-\tprate) = \sigma_t^2$ yields \cref{eq:approx_expected_precision}.

For the variance, the first-order delta method gives $\Var{\PPV} \approx (\nabla f)^\intercal \Sigma\,(\nabla f)$ with $\Sigma = \operatorname{diag}(\sigma_t^2, \sigma_f^2)$:
\begin{equation}
    \Var{\PPV} \approx \frac{\overline{f}^{\,2}\,\sigma_t^2 + \overline{t}^{\,2}\,\sigma_f^2}{\left(\overline{t} + \overline{f}\right)^4},
\end{equation}
which yields \cref{eq:var_ppv}.
\end{proof}

\section{Proof of theorem~\ref{thm:f1_approx}}
\label{sec:proof_f1}
The proof follows the same structure as the PPV proof (\cref{sec:proof_approx_expected_precision}) with the substitution $h(t,f) = 2t/(t + f + \p)$ in place of $f(t,f) = t/(t+f)$.
\begin{proof}
The $\fone$-score is a random variable $\Fone = h(\TP_\p, \FP_\n)$ where
\begin{equation}
    h(t, f) = \frac{2t}{t + f + \p}.
\end{equation}
The gradient of $h$ is
\begin{equation}
    \nabla h(t,f) =
    \frac{2}{(t + f + \p)^2}
    \begin{bmatrix}
        f + \p\\
        -t\\
    \end{bmatrix}
\end{equation}
and the Hessian is
\begin{equation}
    \mathcal{H}(t,f) =
    \frac{2}{(t + f + \p)^3}
    \begin{bmatrix}
        -2(f+\p) & t - f - \p\\
        t - f - \p & 2 t \\
    \end{bmatrix}.
\end{equation}
The second-order Taylor expansion about the means $\overline{t} = \p\tprate$, $\overline{f} = \n\fprate$, letting $s = \overline{t} + \overline{f} + \p$, gives
\begin{equation}
    \Expect{h(\TP_\p, \FP_\n)} \approx h(\overline{t}, \overline{f}) + \frac{1}{2}\operatorname{tr}\!\bigl(\mathcal{H}(\overline{t},\overline{f})\,\Sigma\bigr)
\end{equation}
where $\Sigma = \operatorname{diag}(\sigma_t^2, \sigma_f^2)$ and the cross terms vanish by independence.
Expanding:
\begin{equation}
    \frac{1}{2}\operatorname{tr}\!\bigl(\mathcal{H}\,\Sigma\bigr) =
    \frac{1}{2}\Bigl(\frac{-4(\overline{f}+\p)}{s^3}\,\sigma_t^2 +
    \frac{4\overline{t}}{s^3}\,\sigma_f^2\Bigr) =
    \frac{2\overline{t}\,\sigma_f^2 - 2(\overline{f}+\p)\sigma_t^2}{s^3}.
\end{equation}
Adding $h(\overline{t},\overline{f}) = 2\overline{t}/s$ completes the expected-value formula.

For the variance, the first-order delta method gives $\Var{\Fone} \approx (\nabla h)^\intercal \Sigma\,(\nabla h)$:
\begin{equation}
    \Var{\Fone} \approx \frac{4(\overline{f}+\p)^2 \sigma_t^2 + 4\overline{t}^{\,2}\,\sigma_f^2}{s^4}.
\end{equation}
Substituting $\sigma_t^2 = \p\tprate(1-\tprate)$ and $\sigma_f^2 = \n\fprate(1-\fprate)$ yields \cref{eq:f1_approx,eq:var_f1}.
\end{proof}

\section{Sampling distribution of arbitrary functions}
\label{app:samp}
Suppose we have a function $f : \PS{\Set{X}[1]} \times \cdots \times \PS{\Set{X}[q]} \to \Set{Y}$, and we are interested in evaluating the loss when we replace one or more of the objective input sets with particular corresponding random approximate sets.
The result of this substitution, as previously described, is a probability distribution over $\Set{Y}$.

The probability distribution of random approximate sets are precisely given; therefore, we may estimate the distribution of any function of random approximate sets by generating the random approximate sets and applying the function of interest.

Consider the $m$-by-$q$ matrix where the $(i,j)$-th element is the random approximate set $\tilde{A}_{i\,j}$ such that each random element of the matrix is independently and each column is identically distributed.
If we apply $\Fun{g}$ to each row of the matrix,
\begin{equation}
\RV{Y}_i = \Fun{g}\left(\tilde{A}_{i\,1},\ldots,\tilde{A}_{i\,q}\right)
\end{equation}
for $i=1,\ldots,m$, we generate $m$ i.i.d. random elements $\RV{Y_1},\ldots,\RV{Y_m}$.

If $\Set{Y}$ is a measure space (discrete or continuous), consider the random variable
\begin{equation}
\RV{\overline{Y}_m} = \frac{1}{m} \sum_{i=1}^m \RV{Y_i}.
\end{equation}
If $\RV{Y_1}$ has a well-defined mean $\mu = \Expect{\RV{Y_1}}$ and variance $\sigma^2 = \Var{\RV{Y_1}}$, then by the central limit theorem
\begin{equation}
\sqrt{m}\left(\RV{\overline{Y}_m} - \mu\right) \xrightarrow{d} \normdist\!\left(0, \sigma^2\right)
\end{equation}
as $m \to \infty$.
That is, for large $m$ the sample mean $\RV{\overline{Y}_m}$ is approximately normally distributed with mean $\mu$ and variance $\sigma^2 / m$.

This framework underlies the Monte Carlo validation of the delta method approximations in \cref{fig:ppv_theory_vs_sim,fig:f1_theory_vs_sim}: for each parameter value, $m = 50\,000$ independent random approximate sets are generated, and the sample mean of the measure is compared to the theoretical prediction.
The large sample size ensures that the Monte Carlo standard error is negligible relative to the quantities being estimated.


\bibliography{references}

\end{document}
